\documentclass[12pt]{article}
\usepackage[utf8]{inputenc}
\usepackage{amsmath, amssymb}
\usepackage{algorithm}
\usepackage{algpseudocode}
\usepackage{geometry}
\usepackage{booktabs}
\usepackage{tabularx}
\usepackage{xcolor}
\usepackage{hyperref}
\geometry{a4paper, margin=1in}

\hypersetup{colorlinks=true, linkcolor=blue, urlcolor=blue}

\title{\textbf{Police Patrol Area Covering Models}\\[0.4em]
\large Algorithm Comparison: Curtin et al.\ (2010), Existing Coverage Evaluator,\\
Greedy Optimizer, and Integer Programming Optimizer}
\author{Graduation Project --- Determining Police Patrol Areas}
\date{\today}

\begin{document}
\maketitle
\tableofcontents
\newpage

% =============================================================================
\section{Common Parameters}
% =============================================================================

The following parameters, objectives, and constraints are shared identically
across all four approaches.

\subsection*{Decision Variables}

\begin{tabular}{@{} l p{11cm} @{}}
$I,\,i$            & Set and index of incident locations (calls for service) \\
$J,\,j$            & Set and index of candidate command-centre locations \\
$a_i$              & Weight of incident $i$ \\
$P$                & Number of patrol command centres to locate \\
$S$                & Acceptable service distance threshold \\
$d_{ij}$           & Road-network shortest-path distance, incident $i$ to candidate $j$ \\
$N_i$              & $\{j \in J \mid d_{ij} \le S\}$ — candidates that can cover incident $i$ \\
$x_j \in \{0,1\}$  & 1 if a patrol is located at candidate $j$; 0 otherwise \\
$y_i \in \{0,1\}$  & 1 if incident $i$ is covered by at least one located patrol; 0 otherwise \\
\end{tabular}

\subsection*{Shared Objective and Constraints}

\begin{equation}
    \text{Maximise} \quad O = \sum_{i \in I} a_i\, y_i
    \label{eq:obj}
\end{equation}

\noindent subject to:
\begin{align}
    \sum_{j \in N_i} x_j &\ge y_i  && \forall\, i \in I \label{eq:cov}\\
    \sum_{j \in J}   x_j &=  P     && \label{eq:card}\\
    x_j &\in \{0,1\}               && \forall\, j \in J \label{eq:xbin}\\
    y_i &\in \{0,1\}               && \forall\, i \in I \label{eq:ybin}
\end{align}

\noindent The backup covering objective is additionally computed post-hoc by
all approaches (or formally in Curtin et al.):

\begin{equation}
    B = \sum_{i \in I} a_i \cdot c_i
    \quad \text{where } c_i = |\{j \in H : d_{ij} \le S\}|
    \label{eq:backup}
\end{equation}

\noindent All four approaches use road-network shortest-path distances for
$d_{ij}$, computed via single-source Dijkstra with cutoff $S$ on an OSM
or GIS drivable graph, and evaluate both $O$ and $B$.

% =============================================================================
\section{Differing Parameters}
% =============================================================================

\renewcommand{\arraystretch}{1.45}
\begin{tabularx}{\textwidth}{@{} l X X X X @{}}
\toprule
\textbf{Parameter}
  & \textbf{Curtin et al.\ (2010)}
  & \textbf{Evaluator}
  & \textbf{Greedy}
  & \textbf{IP Optimizer} \\
\midrule

$|I|$
  & Up to 87,603 (Dallas PD, 2000)
  & 997,488 (LAPD)
  & 997,488 (LAPD)
  & 997,488 (LAPD) \\

$|J|$
  & Beat centroids within one division (pre-existing polygons)
  & 21 real stations (fixed; no generation step)
  & 300 K-Means centroids (data-derived)
  & 300 K-Means centroids (data-derived) \\

$P$
  & Per division (e.g.\ 5 for North Central)
  & 21 — all facilities used; no selection
  & 21
  & 21 \\

$S$
  & 2 mi (1 mi for Central Division)
  & 2.5 mi $= 4{,}023$ m
  & 2.5 mi $= 4{,}023$ m
  & 2.5 mi $= 4{,}023$ m \\

$a_i$
  & Signal Code 1–5 (non-linear priority)
  & \texttt{crime\_weight}
  & \texttt{crime\_weight}
  & \texttt{crime\_weight} aggregated to beat level: $a_j = \sum_{i \in \mathcal{S}_j} a_i$ \\

Road network
  & Proprietary NCTCOG GIS (Dallas)
  & OSMnx EPSG:4326
  & OSMnx UTM + EPSG:4326
  & OSMnx UTM + EPSG:4326 \\

Facility selection
  & CPLEX 8.1 IP — proven optimal
  & None (all $J$ used, $x_j \equiv 1$)
  & Greedy $\approx (1-1/e)$ guarantee
  & PuLP + CBC IP — proven optimal \\

IP demand level
  & Incident-level ($|I|$ variables)
  & N/A
  & N/A
  & Beat-level ($|J|=300$ variables) \\

Scope
  & Per division independently (6 runs)
  & City-wide, single pass
  & City-wide, single pass
  & City-wide, single pass \\

Backup coverage
  & Formal multi-objective IP; Pareto frontier of $O$ vs.\ $B$
  & Post-hoc via eq.\ \eqref{eq:backup}
  & Post-hoc via eq.\ \eqref{eq:backup}
  & Post-hoc via eq.\ \eqref{eq:backup} \\

\bottomrule
\end{tabularx}

% =============================================================================
\section{Algorithm 1 — Curtin et al.\ (2010)}
% =============================================================================

\begin{algorithm}[H]
\caption{Classical PPAC --- Curtin, Hayslett-McCall \& Qiu (2010)}
\begin{algorithmic}[1]
\State \textbf{Input:} Incidents $I$ with Signal Code weights $a_i$;
       existing beat-centroid polygons as $J$; $P$; $S$
\State Compute all-to-all OD matrix on proprietary NCTCOG road network (GIS)
\For{each $i \in I$}
    \State $N_i \leftarrow \{j \in J \mid d_{ij} \le S\}$
\EndFor
\State Export $\{|I|,\, |J|,\, a_i,\, N_i\}$ to CPLEX 8.1
\State \textbf{Solve} (branch-and-bound + LP relaxation, proven optimal):
    \Statex \quad $\max \sum_{i} a_i y_i$ \; s.t.\ \eqref{eq:cov}--\eqref{eq:ybin}
\State Assign beats to sectors from $\mathbf{x}^*$; report $O$, total distance
\State \textit{Backup variant:} relax $y_i \in \{0,\dots,P\}$;
       add $\sum_i a_i w_i \ge \hat{O}$;
       solve for decreasing $\hat{O}$ $\to$ Pareto frontier of $(O,B)$
\end{algorithmic}
\end{algorithm}

% =============================================================================
\section{Algorithm 2 — Existing Infrastructure Evaluator}
% =============================================================================

\begin{algorithm}[H]
\caption{Existing LAPD Station Coverage Evaluator}
\begin{algorithmic}[1]
\State \textbf{Input:} Crime CSV ($I$, $a_i$); 21 fixed station locations as $J$;
       city boundary; $S = 2.5$ mi
\State Load, reproject, and spatial-join incidents to city boundary
\State Convert stations from EPSG:2229 $\to$ EPSG:3857; snap all to OSM nodes
\State \textit{No facility selection:} $x_j = 1$ for all $j \in J$
\For{each station $j \in J$}
    \State $\ell_j \leftarrow \text{Dijkstra}(G,\; v^*_j,\; \text{cutoff}=S)$
    \For{each incident $i$ with $v^*_i \in \ell_j$}
        \State $c_i \mathrel{+}= 1$
    \EndFor
\EndFor
\State $O \leftarrow \sum_i a_i \cdot \mathbf{1}[c_i \ge 1]$;\quad
       $B \leftarrow \sum_i a_i \cdot c_i$
\State Report $O$, $B$, coverage percentages; render map
\end{algorithmic}
\end{algorithm}

% =============================================================================
\section{Algorithm 3 — Greedy PPAC Optimizer}
% =============================================================================

\begin{algorithm}[H]
\caption{Greedy PPAC Optimizer}
\begin{algorithmic}[1]
\State \textbf{Input:} Crime CSV ($I$, $a_i$); city boundary; $|J|=300$;
       $P=21$; $S=2.5$ mi
\State \textit{Stage 1 — Candidate generation:}
    Weighted Mini-Batch K-Means on $I$ $\to$ 300 beat centroids $= J$
\State Load OSM graph; project to UTM ($G_{\text{metric}}$) and EPSG:4326 ($G_{\text{4326}}$)
\State Snap all $j \in J$ and $i \in I$ to OSM nodes via $G_{\text{4326}}$ (vectorised)
\For{each candidate $j \in J$}
    \State $C_j \leftarrow \{i \in I \mid v^*_i \in
           \text{Dijkstra}(G_{\text{metric}},\, v^*_j,\, \text{cutoff}=S)\}$
\EndFor
\State \textit{Stage 2 — Greedy selection:}
    $H \leftarrow \emptyset$; \; $\text{covered} \leftarrow \emptyset$
\For{$k = 1$ \textbf{to} $P$}
    \State $j^* \leftarrow \arg\max_{j \notin H}
           \displaystyle\sum_{i \,\in\, C_j \setminus \text{covered}} a_i$
    \State $H \leftarrow H \cup \{j^*\}$;\quad
           $\text{covered} \leftarrow \text{covered} \cup C_{j^*}$
\EndFor
\State $O \leftarrow \sum_i a_i \cdot \mathbf{1}[i \in \text{covered}]$;\quad
       $B \leftarrow \sum_i a_i \cdot |\{j \in H : i \in C_j\}|$
\State Assign beats to nearest $j \in H$; render Voronoi sector map
\end{algorithmic}
\end{algorithm}

% =============================================================================
\section{Algorithm 4 — Integer Programming PPAC Optimizer}
% =============================================================================

\begin{algorithm}[H]
\caption{IP PPAC Optimizer (PuLP + CBC)}
\begin{algorithmic}[1]
\State \textbf{Input:} Crime CSV ($I$, $a_i$); city boundary;
       $|J|=300$; $P=21$; $S=2.5$ mi
\State \textit{Stage 1 — Candidate generation:}
    Weighted Mini-Batch K-Means on $I$ $\to$ $J$;
    aggregate $a_j = \sum_{i \in \mathcal{S}_j} a_i$
\State Load OSM graph; project to $G_{\text{metric}}$ (UTM) and $G_{\text{4326}}$
\State Snap $J$ and $I$ to OSM nodes via $G_{\text{4326}}$ (vectorised)
\For{each candidate $j \in J$}
    \State $C_j \leftarrow \{j' \in J \mid v^*_{j'} \in
           \text{Dijkstra}(G_{\text{metric}},\, v^*_j,\, \text{cutoff}=S)\}$
           \Comment{beat-level coverage sets}
\EndFor
\State \textit{Stage 2 — PPAC IP (PuLP + CBC, proven optimal):}
\State Declare $x_j, y_j \in \{0,1\}$ for all $j \in J$
\State \textbf{Solve:} $\max \sum_{j} a_j y_j$ \; s.t.
    $\sum_{j' \in C_j} x_{j'} \ge y_j \;\forall j$;\;
    $\sum_j x_j = P$;\;
    $x_j, y_j \in \{0,1\}$
\State $H \leftarrow \{j \in J \mid x^*_j = 1\}$
\State \textit{Stage 3 — Incident-level coverage evaluation:}
\For{each $j \in H$}
    \State $C^{\text{inc}}_j \leftarrow \{i \in I \mid v^*_i \in
           \text{Dijkstra}(G_{\text{metric}},\, v^*_j,\, \text{cutoff}=S)\}$
    \For{each $i \in C^{\text{inc}}_j$}
        \State $c_i \mathrel{+}= 1$
    \EndFor
\EndFor
\State $O \leftarrow \sum_i a_i \cdot \mathbf{1}[c_i \ge 1]$;\quad
       $B \leftarrow \sum_i a_i \cdot c_i$
\State Assign beats to nearest $j \in H$; render Voronoi sector map
\end{algorithmic}
\end{algorithm}

% =============================================================================
\begin{thebibliography}{9}
\bibitem{curtin2010}
Curtin, K.M., Hayslett-McCall, K., \& Qiu, F. (2010).
\textit{Determining Optimal Police Patrol Areas with Maximal Covering and
Backup Covering Location Models.}
Networks and Spatial Economics, 10, 125--145.

\bibitem{church1974}
Church, R.L., \& ReVelle, C.S. (1974).
\textit{The Maximal Covering Location Problem.}
Papers of the Regional Science Association, 32, 101--118.
\end{thebibliography}

\end{document}